When the unit ring is local, the ordinary projective data in this
reduction have a complete multiplicity description. Here a local ring
is nonzero and may be noncommutative.

\begin{theorem}[local unit rings]
\label{thm:local-unit-projectives}
Suppose $T=(Q,F)$ has integral support and $R=F(1)$ is local. Put
$I_q=\im(F(q)\to R)$ and $W=\{q:I_q=R\}$. For a finitely generated
projective left $Q$-support $L$, define
\[
 l\preceq_W k\quad\Longleftrightarrow\quad
 ql\leq k\text{ for some }q\in W,
\]
and let $l\sim_W k$ mean comparability in both directions. A class
$[l]$ is \emph{allowed} if a functional $\psi$ admissible at $l$
satisfies $\psi(l)\in W$, equivalently if $C(l,l)\cap W\ne\varnothing$.
This condition is independent of its representative.

Finite-projective left $T$-semimodules are classified by $L$ and a
finitely supported function
\[
 m:\{\text{allowed }\sim_W\text{ classes in }L\}\longrightarrow\N,
\]
up to $Q$-support isomorphisms transporting these multiplicities.
For the separated diagram of Theorem~\ref{thm:integral-support-reduction},
choose representatives of the classes with positive multiplicity. Its
group over $k$ is
\begin{equation}\label{eq:local-unit-fibers}
 G_k=\bigoplus_{[l]}
 \left(\bigcup_{ql\leq k}I_q\right)^{m([l])}.
\end{equation}
Transitions are inclusions and scalars act by left multiplication in
$R$. The original $T$-class is recovered by lifting an idempotent; it
is not obtained by an assertion that restriction of scalars preserves
projectivity.
\end{theorem}

The theorem evaluates the ordinary amplitude classification for every
integral support with local unit ring, including infinite supports and
infinite coefficient groups. Classifying the projective $Q$-supports
themselves remains a separate problem. Zero multiplicities are allowed;
the resulting zero-amplitude module still has its specified support.
